\chapter{\nt{ГЛУТ} и \nt{ГАМК} вместе}
\label{sec:gaba}

\section{Правила игры}

Сеть из одних \nt{ГЛУТ}-нейронов не умеет выбирать, сбрасывать память и удерживать активность от лавины, и всё из-за монотонности (утверждение~\ref{prop:monotone}). Добавим в неё второй по численности тип нейронов --- \nt{ГАМК}. Правила прежние: только то, что бывает в живом организме, и никаких часов.

Модель нейрона та же, что~\eqref{eq:glutneuron}, но импульс от \nt{ГАМК}-нейрона не поднимает напряжение получателя, а опускает. Веса теперь бывают двух знаков, и знак задаёт отправитель, правило~\eqref{eq:dalesign}.

Несколько параметров схемы. \nt{ГАМК}-нейронов в коре от пятой части до четверти~\cite{Hendry1987}. Выходы у них короткие: они тормозят соседей, а не далёкие области. Торможение приходит через миллисекунду и длится несколько миллисекунд. Без входа \nt{ГАМК}-нейрон молчит: чтобы затормозить кого-то, он сам должен получить возбуждение, обычно от \nt{ГЛУТ}-нейронов той же сети.

Поэтому отрицательную связь от \nt{ГЛУТ}-нейрона $A$ к нейрону $B$ приходится делать через посредника: $A$ возбуждает \nt{ГАМК}-нейрон, а тот тормозит $B$. Смена знака стоит одного нейрона и одной задержки, около миллисекунды.

Для торможения важно не только, \emph{от кого} идёт связь, но и \emph{куда} она садится: место посадки во многом определяет, что связь делает~\cite{Markram2004}. \nt{ГЛУТ}-выходы садятся в основном на входные ветви получателя, его \emph{дендриты}, и несут данные: каждый такой вход --- одно слагаемое суммы. Выход на тело нейрона действует на сумму целиком: так многие быстрые \nt{ГАМК}-нейроны делят ответ получателя и задают, когда ему сработать. Выход на то место, где рождается импульс получателя, --- последнее слово, запрет на весь выход сразу. А выход на окончание чужого провода меняет вероятность выброса $p$ одной входной линии, не трогая остальных~\cite{Rudomin1999}; так, в частности, работают ворота боли в главе~\ref{sec:pain}. За пределами мозга выходы садятся на мышечные волокна (глава~\ref{sec:muscles}) и на органы (глава~\ref{sec:chemoutput}), а некоторые не садятся никуда и выбрасывают медиатор в объём ткани или в кровь (главы~\ref{sec:serocomputer} и~\ref{sec:chemoutput}).

\section{НЕ и запрет}

Можно ли теперь сделать НЕ? Почти. Нейрон, который работает при молчащем входе, должен откуда-то брать возбуждение. В живом мозге оно есть: на каждый нейрон постоянно приходит фоновая активность от тысяч других. Пусть фон держит нейрон $Y$ в работе, а вход $x$ через \nt{ГАМК}-посредника его тормозит. Это НЕ.

Полезнее, однако, \emph{запрет}: «$a$, но не при $b$»:
\begin{empheq}[box=\keybox]{equation}
y \;=\; a \wedge \bar b .
\label{eq:veto}
\end{empheq}
Нейрон $Y$ получает возбуждение от $a$ и торможение от \nt{ГАМК}-нейрона, которого возбуждает $b$. Фон здесь не нужен: возбуждение даёт сам $a$. Из запретов и ИЛИ собирается всё. Например, исключающее ИЛИ: $x\oplus y=(x\wedge\bar y)\vee(\bar x\wedge y)$ --- два нейрона-запрета, два \nt{ГАМК}-посредника и один нейрон ИЛИ.

\begin{keyblock}
\begin{proposition}[Полнота \nt{ГЛУТ}+\nt{ГАМК}]
\label{prop:gabacomplete}
Сеть из \nt{ГЛУТ}- и \nt{ГАМК}-нейронов может вычислить любую логическую функцию входов, и каждый бит передаётся одним проводом. Датчики «включения---выключения» из главы~\ref{sec:glutcomputer} для этого больше не нужны. Если функция должна давать единицу при всех молчащих входах, нужен источник фоновой активности.
\end{proposition}
\end{keyblock}

Доказательство: запрет~\eqref{eq:veto} вместе с ИЛИ функционально полон, если есть постоянная единица, потому что НЕ $x$ --- это запрет «единица, но не при $x$». А функции, которые при всех молчащих входах дают ноль, собираются из запретов и ИЛИ и без единицы.

\section{Направление движения}

Запрет во времени, как и И во времени в главе~\ref{sec:glutcomputer}, даёт детектор направления движения.

Пусть два соседних датчика, $A$ и $B$, стоят на расстоянии $\Delta x$ друг от друга. Выходной нейрон $Y$ получает возбуждение от $B$ и торможение от $A$ через \nt{ГАМК}-посредника, с задержкой $d$ (рис.~\ref{fig:direction}). Пятно света, бегущее со скоростью $v$ от $A$ к $B$, проходит $A$ в момент $0$, и торможение приходит к $Y$ в момент $d$; $B$ оно проходит в момент $\Delta x/v$, и тогда же приходит возбуждение. Если эти моменты совпадают с точностью до окна~\eqref{eq:window}, торможение гасит возбуждение, и $Y$ молчит. Это происходит при скорости около
\begin{empheq}[box=\keybox]{equation}
v^* \;=\; \frac{\Delta x}{d} .
\label{eq:vstar}
\end{empheq}
При движении от $B$ к $A$ возбуждение от $B$ приходит в момент $0$, а торможение от $A$ --- только в момент $\Delta x/v+d$, когда $Y$ уже сработал.

\begin{figure}[t]
\centering
\begin{tikzpicture}[>=Stealth,
  nrn/.style={circle,draw,very thick,fill=blue!6,minimum size=0.9cm,inner sep=0pt},
  gab/.style={nrn,fill=red!10},
  inh/.style={-{Bar[width=7pt]},thick,red!70!black}]
\node[nrn] (A) at (0,2) {$A$};
\node[nrn] (B) at (3,2) {$B$};
\node[gab] (G) at (0.9,0.6) {\small \nt{ГАМК}};
\node[nrn] (Y) at (3,-0.6) {$Y$};
\draw[->,thick] (A) -- (G);
\draw[inh] (G) -- node[below left=-2pt] {\scriptsize задержка $d$} (Y);
\draw[->,thick] (B) -- (Y);
\draw[->,very thick,red!70!black] (Y) -- (4.6,-0.6) node[right,black] {\small выход};
\draw[->,thick,orange!85!black] (-0.6,2.9) -- (3.6,2.9) node[right,black] {\small молчит};
\draw[<-,thick,orange!85!black] (-0.6,3.4) -- (3.6,3.4) node[right,black] {\small отвечает};
\foreach \yy/\dd in {2.9/-1,3.4/1}{
  \foreach \k in {1,2,3}{\draw[orange!70!yellow,line width=0.8pt] (1.5+\dd*0.12+\dd*0.13*\k,\yy+0.1-0.1*\k+0.1) -- ++(\dd*0.25,0);}
  \fill[yellow!70!orange,draw=orange!70!black] (1.5,\yy) circle (0.14);}
\node[font=\scriptsize,align=center] at (1.5,3.95) {пятно света};
\end{tikzpicture}
\caption{Детектор направления движения. $B$ возбуждает $Y$, $A$ тормозит его через \nt{ГАМК}-посредника с задержкой $d$. При движении от $A$ к $B$ со скоростью около $\Delta x/d$ торможение гасит возбуждение; при обратном движении оно опаздывает. Жёлтое пятно с чёрточками --- свет, бегущий над датчиками.}
\label{fig:direction}
\end{figure}

В сетчатке избирательность к направлению движения, по-видимому, создаётся именно так: асимметричным торможением со стороны, откуда движение не должно вызывать ответ~\cite{Barlow1965}.

\section{Вычитание или деление}

До сих пор торможение у нас вычитало. На самом деле оно устроено тоньше.

Синапс открывает канал, то есть включает проводимость. У возбуждающего синапса равновесный потенциал лежит далеко выше порога, примерно на $70\mV$ выше покоя: открытый канал тянет напряжение вверх. У тормозящего \nt{ГАМК}-синапса канал пропускает хлор, равновесный потенциал которого лежит около покоя, и открытый канал тянет напряжение не вниз, а \emph{к покою}. Если отсчитывать напряжение от покоя, установившееся напряжение мембраны с проводимостью утечки $g_L$, возбуждающей $g_E$ и тормозящей $g_I$ равно
\begin{empheq}[box=\keybox]{equation}
V_\infty \;=\; \frac{g_E\,E_E}{g_L+g_E+g_I} ,
\label{eq:shunt}
\end{empheq}
где $E_E\approx70\mV$ --- равновесный потенциал возбуждающего синапса. Это закон Ома для делителя напряжения: $g_E$ тянет к $E_E$, а $g_L$ и $g_I$ --- к нулю.

$g_I$ стоит не в числителе со знаком минус, а в знаменателе: торможение не \emph{вычитает} из возбуждения, а \emph{делит} его (рис.~\ref{fig:shunt}). Такое торможение называют шунтирующим: открытые хлорные каналы закорачивают мембрану. Для сети это регулировка усиления. Если $g_I$ пропорциональна суммарной активности соседей, каждый нейрон отвечает не на абсолютную силу своего входа, а на её долю в общей активности. Такое деление на общую активность встречается во многих отделах мозга и считается одной из его стандартных операций~\cite{Carandini2012}: одна и та же сеть работает и при слабом, и при сильном входе.

Оговорка: формула~\eqref{eq:shunt} --- про нейрон, который не выдаёт импульсов. У срабатывающего нейрона напряжение всё время держится около порога, ток через тормозящую проводимость почти постоянен, и на \emph{частоту} импульсов шунт действует в основном вычитанием~\cite{HoltKoch1997}. Частота делится, если нейрону одновременно добавить уравновешенные возбуждающую и тормозящую фоновые проводимости, как в дрожащем, уравновешенном режиме живой коры~\cite{Chance2002}. Так что деление мозг получает не от одного шунта, а от торможения вместе с фоновым шумом~\cite{Carandini2012}.

\begin{figure}[t]
\centering
\begin{tikzpicture}[>=Stealth,x=1.25cm,y=0.05cm]
\draw[->,gray!70!black] (0,0) -- (5.4,0) node[right,black] {\small $g_E/g_L$};
\draw[->,gray!70!black] (0,0) -- (0,68) node[above,black] {\small $V_\infty$, мВ};
\foreach \v in {20,40,60}{\draw[gray!60!black] (-0.05,\v) -- (0.05,\v) node[left=3pt,font=\scriptsize] {\v};}
\foreach \x in {1,2,3,4,5}{\draw[gray!60!black] (\x,-1.5) -- (\x,1.5) node[below=5pt,font=\scriptsize] {\x};}
\draw[very thick,blue!60!black] plot[domain=0:5,samples=60] (\x,{70*\x/(1+\x)});
\draw[very thick,red!70!black] plot[domain=0:5,samples=60] (\x,{70*\x/(3+\x)});
\draw[thick,densely dashed,gray!50!black] plot[domain=0.56:5,samples=60] (\x,{70*\x/(1+\x)-25});
\node[right,font=\small,blue!60!black] at (5.02,58.3) {без торможения};
\node[right,font=\small,red!70!black] at (5.02,45.5) {делит: $g_I=2g_L$};
\node[right,font=\small,gray!40!black] at (5.02,32.5) {вычитает $25\mV$};
\end{tikzpicture}
\caption{Шунтирующее торможение делит напряжение, а не вычитает из него. Синяя кривая --- установившееся напряжение~\eqref{eq:shunt} без торможения, красная --- с тормозящей проводимостью $g_I=2g_L$. Красная --- та же синяя, растянутая по горизонтали втрое: вход как будто поделили на $1+g_I/g_L=3$. Штриховая --- торможение, вычитающее постоянные $25\mV$: кривая опускается целиком, наклон не меняется.}
\label{fig:shunt}
\end{figure}

Есть и второе следствие. Постоянная времени мембраны $\tau_{\text{эфф}}=C/(g_L+g_E+g_I)$, и торможение её укорачивает. А окно совпадения~\eqref{eq:window} пропорционально $\tau$, так что торможение, пришедшее вместе с возбуждением, \emph{сужает окно}. Схема, в которой вход возбуждает нейрон напрямую и одновременно, с задержкой в миллисекунду, тормозит его через \nt{ГАМК}-посредника, --- одна из самых распространённых в мозге: нейрон успевает ответить только на то, что пришло в первую миллисекунду-другую~\cite{Pouille2001}.

\section{Выбор}

Теперь главное, чего не хватало \nt{ГЛУТ}-сети: выбор одного из нескольких. Возьмём две группы \nt{ГЛУТ}-нейронов с входами $I_1$ и $I_2$. Каждая через своих \nt{ГАМК}-посредников тормозит другую с силой $\beta$ (рис.~\ref{fig:wta}). Для частот $r_1,r_2$ получаем
\begin{equation}
\tau\,\frac{\dd r_1}{\dd t} = -r_1 + \bigl[I_1-\beta r_2\bigr]_+ ,\qquad
\tau\,\frac{\dd r_2}{\dd t} = -r_2 + \bigl[I_2-\beta r_1\bigr]_+ ,
\label{eq:wta}
\end{equation}
где $[u]_+=\max(u,0)$: частота не бывает отрицательной.

\begin{figure}[t]
\centering
\begin{tikzpicture}[>=Stealth,
  nrn/.style={circle,draw,very thick,fill=blue!6,minimum size=0.9cm,inner sep=0pt},
  gab/.style={nrn,fill=red!10,minimum size=0.75cm},
  inh/.style={-{Bar[width=7pt]},thick,red!70!black}]
\node[nrn] (E1) at (0,0) {$r_1$};
\node[nrn] (E2) at (4,0) {$r_2$};
\node[gab] (G1) at (1.4,1.1) {};
\node[gab] (G2) at (2.6,-1.1) {};
\draw[->,thick] (E1) -- (G1); \draw[inh] (G1) -- (E2);
\draw[->,thick] (E2) -- (G2); \draw[inh] (G2) -- (E1);
\draw[->,thick,blue!60!black] (-1.6,0) node[left,black] {$I_1$} -- (E1);
\draw[->,thick,blue!60!black] (5.6,0) node[right,black] {$I_2$} -- (E2);
\end{tikzpicture}
\caption{Выбор из двух. Две группы \nt{ГЛУТ}-нейронов (синие) тормозят друг друга через \nt{ГАМК}-посредников (красные).}
\label{fig:wta}
\end{figure}

Пока работают оба нейрона, система~\eqref{eq:wta} линейна, и её поведение задают два собственных числа: $-(1+\beta)/\tau$ для одинаковых отклонений обеих частот и $-(1-\beta)/\tau$ для противоположных. При слабом торможении, $\beta<1$, оба числа отрицательны: оба нейрона работают, торможение лишь приглушает слабого. При сильном торможении, $\beta>1$, второе число становится положительным: любая разница между $r_1$ и $r_2$ нарастает, пока один нейрон не замолчит совсем.

\begin{keyblock}
\begin{proposition}[Победитель забирает всё]
\label{prop:wta}
Если взаимное торможение сильнее единицы, $\beta>1$, состояние, в котором работают обе группы, неустойчиво. Сеть приходит в состояние, где работает одна группа, а другая молчит. Победитель с частотой $r_1=I_1$ удерживает победу, пока $I_2<\beta I_1$.
\end{proposition}
\end{keyblock}

Мы проверили это на модели. При $\beta=0.5$ и входах $1.0$ и $0.9$ обе группы работают, с частотами $0.73$ и $0.53$. При $\beta=2$ и тех же входах вторая группа молчит совсем. У такого выбора есть память: если потом поменять входы местами, прежний победитель остаётся победителем, потому что $1.0<2\cdot0.9$.

С сотней групп всё так же: каждая тормозит остальных, выживает одна.

\section{Сброс памяти}

Память из главы~\ref{sec:glutcomputer} выключалась только усталостью. Добавим к ней \nt{ГАМК}-нейрон, который возбуждается сигналом «сброс» и тормозит группу. Торможение опускает кривую $f(wr+I)$ на рис.~\ref{fig:bistable}, верхнее пересечение исчезает, и группа гаснет --- и остаётся в нижнем состоянии, когда сброс кончился. Теперь память записывается одним сигналом и стирается другим, как у триггера со входами установки и сброса.

Ещё красивее соединить две такие группы взаимным торможением, как на рис.~\ref{fig:wta}, и дать каждой обратную связь на саму себя. Сигнал на вход одной из групп переводит ячейку в её состояние, и ячейка помнит последнее решение. Это прямой аналог защёлки из двух перекрёстно соединённых элементов ИЛИ-НЕ.

\section{Устойчивость и ритм}

Остаётся третья беда \nt{ГЛУТ}-сети --- лавина. Возьмём группу \nt{ГЛУТ}-нейронов с обратной связью на себя $w_{EE}$ и группу \nt{ГАМК}-нейронов, которых первая возбуждает с силой $w_{IE}$ и которые тормозят первую с силой $w_{EI}$. Для частот $E$ и $I$ получаем
\begin{equation}
\tau_E\frac{\dd E}{\dd t} = -E + w_{EE}E - w_{EI}I + h, \qquad
\tau_I\frac{\dd I}{\dd t} = -I + w_{IE}E ,
\label{eq:ei}
\end{equation}
где $h$ --- внешний вход~\cite{WilsonCowan1972}. Без торможения при $w_{EE}>1$ активность растёт без предела: каждый импульс порождает больше одного следующего. С торможением установившаяся частота
\begin{equation}
E^* \;=\; \frac{h}{1-w_{EE}+w_{EI}w_{IE}}
\label{eq:eistar}
\end{equation}
конечна, если знаменатель положителен. Но этого мало: равновесие должно ещё и притягивать. Из линейного анализа~\eqref{eq:ei} получаются два условия.

\begin{keyblock}
\begin{proposition}[Условия устойчивости]
\label{prop:ei}
Равновесие~\eqref{eq:eistar} устойчиво, если торможение
\begin{enumerate}
\item[(i)] \emph{достаточно сильное}: $w_{EI}\,w_{IE} > w_{EE}-1$;
\item[(ii)] \emph{достаточно быстрое}: $\tau_I < \tau_E/(w_{EE}-1)$.
\end{enumerate}
Если нарушено (i), активность нарастает лавиной. Если выполнено (i), но нарушено (ii), торможение запаздывает, и активность начинает колебаться.
\end{proposition}
\end{keyblock}

Условие~(i) --- это определитель матрицы системы~\eqref{eq:ei}, условие~(ii) --- её след. Смысл второго понятен любому, кто настраивал регулятор: запоздавшая обратная связь не гасит отклонение, а раскачивает его.

\begin{figure}[t]
\centering
\begin{tikzpicture}[>=Stealth]
\draw[->,gray!70!black] (0,0) -- (10.9,0) node[below left,black] {\small $t$, мс};
\draw[->,gray!70!black] (0,0) -- (0,3.3) node[left,black] {\small $E$};
\foreach \t in {0,50,100,150}{\draw[gray!70!black] (\t*0.07,0.06) -- (\t*0.07,-0.06) node[below,font=\scriptsize] {\t};}
\draw[very thick,gray!60!black,densely dashed] plot coordinates {(0.000,0.000) (0.035,0.071) (0.070,0.144) (0.105,0.218) (0.140,0.294) (0.175,0.373) (0.210,0.453) (0.245,0.535) (0.280,0.620) (0.315,0.706) (0.350,0.795) (0.385,0.886) (0.420,0.979) (0.455,1.075) (0.490,1.173) (0.525,1.274) (0.560,1.377) (0.595,1.482) (0.630,1.591) (0.665,1.702) (0.700,1.816) (0.735,1.933) (0.770,2.052) (0.805,2.175) (0.840,2.301) (0.875,2.430) (0.910,2.563) (0.945,2.698) (0.980,2.838) (1.015,2.980) (1.050,3.080) (1.085,3.080) (1.120,3.080) (1.155,3.080) (1.190,3.080) (1.225,3.080) (1.260,3.080) (1.295,3.080) (1.330,3.080) (1.365,3.080) (1.400,3.080) (1.435,3.080) (1.470,3.080) (1.505,3.080) (1.540,3.080) (1.575,3.080) (1.610,3.080) (1.645,3.080) (1.680,3.080) (1.715,3.080) (1.750,3.080) (1.785,3.080) (1.820,3.080) (1.855,3.080) (1.890,3.080) (1.925,3.080) (1.960,3.080) (1.995,3.080) (2.030,3.080) (2.065,3.080) (2.100,3.080) (2.135,3.080) (2.170,3.080) (2.205,3.080) (2.240,3.080) (2.275,3.080) (2.310,3.080) (2.345,3.080) (2.380,3.080) (2.415,3.080) (2.450,3.080) (2.485,3.080) (2.520,3.080) (2.555,3.080) (2.590,3.080) (2.625,3.080) (2.660,3.080) (2.695,3.080) (2.730,3.080) (2.765,3.080) (2.800,3.080) (2.835,3.080) (2.870,3.080) (2.905,3.080) (2.940,3.080) (2.975,3.080) (3.010,3.080) (3.045,3.080) (3.080,3.080) (3.115,3.080) (3.150,3.080) (3.185,3.080) (3.220,3.080) (3.255,3.080) (3.290,3.080) (3.325,3.080) (3.360,3.080) (3.395,3.080) (3.430,3.080) (3.465,3.080) (3.500,3.080) (3.535,3.080) (3.570,3.080) (3.605,3.080) (3.640,3.080) (3.675,3.080) (3.710,3.080) (3.745,3.080) (3.780,3.080) (3.815,3.080) (3.850,3.080) (3.885,3.080) (3.920,3.080) (3.955,3.080) (3.990,3.080) (4.025,3.080) (4.060,3.080) (4.095,3.080) (4.130,3.080) (4.165,3.080) (4.200,3.080) (4.235,3.080) (4.270,3.080) (4.305,3.080) (4.340,3.080) (4.375,3.080) (4.410,3.080) (4.445,3.080) (4.480,3.080) (4.515,3.080) (4.550,3.080) (4.585,3.080) (4.620,3.080) (4.655,3.080) (4.690,3.080) (4.725,3.080) (4.760,3.080) (4.795,3.080) (4.830,3.080) (4.865,3.080) (4.900,3.080) (4.935,3.080) (4.970,3.080) (5.005,3.080) (5.040,3.080) (5.075,3.080) (5.110,3.080) (5.145,3.080) (5.180,3.080) (5.215,3.080) (5.250,3.080) (5.285,3.080) (5.320,3.080) (5.355,3.080) (5.390,3.080) (5.425,3.080) (5.460,3.080) (5.495,3.080) (5.530,3.080) (5.565,3.080) (5.600,3.080) (5.635,3.080) (5.670,3.080) (5.705,3.080) (5.740,3.080) (5.775,3.080) (5.810,3.080) (5.845,3.080) (5.880,3.080) (5.915,3.080) (5.950,3.080) (5.985,3.080) (6.020,3.080) (6.055,3.080) (6.090,3.080) (6.125,3.080) (6.160,3.080) (6.195,3.080) (6.230,3.080) (6.265,3.080) (6.300,3.080) (6.335,3.080) (6.370,3.080) (6.405,3.080) (6.440,3.080) (6.475,3.080) (6.510,3.080) (6.545,3.080) (6.580,3.080) (6.615,3.080) (6.650,3.080) (6.685,3.080) (6.720,3.080) (6.755,3.080) (6.790,3.080) (6.825,3.080) (6.860,3.080) (6.895,3.080) (6.930,3.080) (6.965,3.080) (7.000,3.080) (7.035,3.080) (7.070,3.080) (7.105,3.080) (7.140,3.080) (7.175,3.080) (7.210,3.080) (7.245,3.080) (7.280,3.080) (7.315,3.080) (7.350,3.080) (7.385,3.080) (7.420,3.080) (7.455,3.080) (7.490,3.080) (7.525,3.080) (7.560,3.080) (7.595,3.080) (7.630,3.080) (7.665,3.080) (7.700,3.080) (7.735,3.080) (7.770,3.080) (7.805,3.080) (7.840,3.080) (7.875,3.080) (7.910,3.080) (7.945,3.080) (7.980,3.080) (8.015,3.080) (8.050,3.080) (8.085,3.080) (8.120,3.080) (8.155,3.080) (8.190,3.080) (8.225,3.080) (8.260,3.080) (8.295,3.080) (8.330,3.080) (8.365,3.080) (8.400,3.080) (8.435,3.080) (8.470,3.080) (8.505,3.080) (8.540,3.080) (8.575,3.080) (8.610,3.080) (8.645,3.080) (8.680,3.080) (8.715,3.080) (8.750,3.080) (8.785,3.080) (8.820,3.080) (8.855,3.080) (8.890,3.080) (8.925,3.080) (8.960,3.080) (8.995,3.080) (9.030,3.080) (9.065,3.080) (9.100,3.080) (9.135,3.080) (9.170,3.080) (9.205,3.080) (9.240,3.080) (9.275,3.080) (9.310,3.080) (9.345,3.080) (9.380,3.080) (9.415,3.080) (9.450,3.080) (9.485,3.080) (9.520,3.080) (9.555,3.080) (9.590,3.080) (9.625,3.080) (9.660,3.080) (9.695,3.080) (9.730,3.080) (9.765,3.080) (9.800,3.080) (9.835,3.080) (9.870,3.080) (9.905,3.080) (9.940,3.080) (9.975,3.080) (10.010,3.080) (10.045,3.080) (10.080,3.080) (10.115,3.080) (10.150,3.080) (10.185,3.080) (10.220,3.080) (10.255,3.080) (10.290,3.080) (10.325,3.080) (10.360,3.080) (10.395,3.080) (10.430,3.080) (10.465,3.080) (10.500,3.080)};
\draw[very thick,blue!60!black] plot coordinates {(0.000,0.000) (0.035,0.071) (0.070,0.141) (0.105,0.211) (0.140,0.279) (0.175,0.344) (0.210,0.406) (0.245,0.465) (0.280,0.519) (0.315,0.570) (0.350,0.616) (0.385,0.658) (0.420,0.697) (0.455,0.731) (0.490,0.763) (0.525,0.790) (0.560,0.815) (0.595,0.836) (0.630,0.855) (0.665,0.871) (0.700,0.886) (0.735,0.898) (0.770,0.908) (0.805,0.916) (0.840,0.924) (0.875,0.929) (0.910,0.934) (0.945,0.938) (0.980,0.941) (1.015,0.943) (1.050,0.945) (1.085,0.946) (1.120,0.946) (1.155,0.947) (1.190,0.947) (1.225,0.947) (1.260,0.946) (1.295,0.946) (1.330,0.945) (1.365,0.944) (1.400,0.944) (1.435,0.943) (1.470,0.942) (1.505,0.941) (1.540,0.941) (1.575,0.940) (1.610,0.939) (1.645,0.939) (1.680,0.938) (1.715,0.937) (1.750,0.937) (1.785,0.936) (1.820,0.936) (1.855,0.936) (1.890,0.935) (1.925,0.935) (1.960,0.935) (1.995,0.934) (2.030,0.934) (2.065,0.934) (2.100,0.934) (2.135,0.934) (2.170,0.934) (2.205,0.934) (2.240,0.933) (2.275,0.933) (2.310,0.933) (2.345,0.933) (2.380,0.933) (2.415,0.933) (2.450,0.933) (2.485,0.933) (2.520,0.933) (2.555,0.933) (2.590,0.933) (2.625,0.933) (2.660,0.933) (2.695,0.933) (2.730,0.933) (2.765,0.933) (2.800,0.933) (2.835,0.933) (2.870,0.933) (2.905,0.933) (2.940,0.933) (2.975,0.933) (3.010,0.933) (3.045,0.933) (3.080,0.933) (3.115,0.933) (3.150,0.933) (3.185,0.933) (3.220,0.933) (3.255,0.933) (3.290,0.933) (3.325,0.933) (3.360,0.933) (3.395,0.933) (3.430,0.933) (3.465,0.933) (3.500,0.933) (3.535,0.933) (3.570,0.933) (3.605,0.933) (3.640,0.933) (3.675,0.933) (3.710,0.933) (3.745,0.933) (3.780,0.933) (3.815,0.933) (3.850,0.933) (3.885,0.933) (3.920,0.933) (3.955,0.933) (3.990,0.933) (4.025,0.933) (4.060,0.933) (4.095,0.933) (4.130,0.933) (4.165,0.933) (4.200,0.933) (4.235,0.933) (4.270,0.933) (4.305,0.933) (4.340,0.933) (4.375,0.933) (4.410,0.933) (4.445,0.933) (4.480,0.933) (4.515,0.933) (4.550,0.933) (4.585,0.933) (4.620,0.933) (4.655,0.933) (4.690,0.933) (4.725,0.933) (4.760,0.933) (4.795,0.933) (4.830,0.933) (4.865,0.933) (4.900,0.933) (4.935,0.933) (4.970,0.933) (5.005,0.933) (5.040,0.933) (5.075,0.933) (5.110,0.933) (5.145,0.933) (5.180,0.933) (5.215,0.933) (5.250,0.933) (5.285,0.933) (5.320,0.933) (5.355,0.933) (5.390,0.933) (5.425,0.933) (5.460,0.933) (5.495,0.933) (5.530,0.933) (5.565,0.933) (5.600,0.933) (5.635,0.933) (5.670,0.933) (5.705,0.933) (5.740,0.933) (5.775,0.933) (5.810,0.933) (5.845,0.933) (5.880,0.933) (5.915,0.933) (5.950,0.933) (5.985,0.933) (6.020,0.933) (6.055,0.933) (6.090,0.933) (6.125,0.933) (6.160,0.933) (6.195,0.933) (6.230,0.933) (6.265,0.933) (6.300,0.933) (6.335,0.933) (6.370,0.933) (6.405,0.933) (6.440,0.933) (6.475,0.933) (6.510,0.933) (6.545,0.933) (6.580,0.933) (6.615,0.933) (6.650,0.933) (6.685,0.933) (6.720,0.933) (6.755,0.933) (6.790,0.933) (6.825,0.933) (6.860,0.933) (6.895,0.933) (6.930,0.933) (6.965,0.933) (7.000,0.933) (7.035,0.933) (7.070,0.933) (7.105,0.933) (7.140,0.933) (7.175,0.933) (7.210,0.933) (7.245,0.933) (7.280,0.933) (7.315,0.933) (7.350,0.933) (7.385,0.933) (7.420,0.933) (7.455,0.933) (7.490,0.933) (7.525,0.933) (7.560,0.933) (7.595,0.933) (7.630,0.933) (7.665,0.933) (7.700,0.933) (7.735,0.933) (7.770,0.933) (7.805,0.933) (7.840,0.933) (7.875,0.933) (7.910,0.933) (7.945,0.933) (7.980,0.933) (8.015,0.933) (8.050,0.933) (8.085,0.933) (8.120,0.933) (8.155,0.933) (8.190,0.933) (8.225,0.933) (8.260,0.933) (8.295,0.933) (8.330,0.933) (8.365,0.933) (8.400,0.933) (8.435,0.933) (8.470,0.933) (8.505,0.933) (8.540,0.933) (8.575,0.933) (8.610,0.933) (8.645,0.933) (8.680,0.933) (8.715,0.933) (8.750,0.933) (8.785,0.933) (8.820,0.933) (8.855,0.933) (8.890,0.933) (8.925,0.933) (8.960,0.933) (8.995,0.933) (9.030,0.933) (9.065,0.933) (9.100,0.933) (9.135,0.933) (9.170,0.933) (9.205,0.933) (9.240,0.933) (9.275,0.933) (9.310,0.933) (9.345,0.933) (9.380,0.933) (9.415,0.933) (9.450,0.933) (9.485,0.933) (9.520,0.933) (9.555,0.933) (9.590,0.933) (9.625,0.933) (9.660,0.933) (9.695,0.933) (9.730,0.933) (9.765,0.933) (9.800,0.933) (9.835,0.933) (9.870,0.933) (9.905,0.933) (9.940,0.933) (9.975,0.933) (10.010,0.933) (10.045,0.933) (10.080,0.933) (10.115,0.933) (10.150,0.933) (10.185,0.933) (10.220,0.933) (10.255,0.933) (10.290,0.933) (10.325,0.933) (10.360,0.933) (10.395,0.933) (10.430,0.933) (10.465,0.933) (10.500,0.933)};
\draw[very thick,red!70!black] plot coordinates {(0.000,0.000) (0.035,0.071) (0.070,0.143) (0.105,0.217) (0.140,0.293) (0.175,0.370) (0.210,0.449) (0.245,0.528) (0.280,0.609) (0.315,0.691) (0.350,0.774) (0.385,0.858) (0.420,0.943) (0.455,1.028) (0.490,1.114) (0.525,1.200) (0.560,1.287) (0.595,1.374) (0.630,1.462) (0.665,1.549) (0.700,1.636) (0.735,1.724) (0.770,1.810) (0.805,1.897) (0.840,1.983) (0.875,2.068) (0.910,2.153) (0.945,2.237) (0.980,2.320) (1.015,2.402) (1.050,2.483) (1.085,2.563) (1.120,2.641) (1.155,2.717) (1.190,2.792) (1.225,2.866) (1.260,2.937) (1.295,3.007) (1.330,3.075) (1.365,3.080) (1.400,3.080) (1.435,3.080) (1.470,3.080) (1.505,3.080) (1.540,3.080) (1.575,3.080) (1.610,3.080) (1.645,3.080) (1.680,3.080) (1.715,3.080) (1.750,3.080) (1.785,3.080) (1.820,3.080) (1.855,3.080) (1.890,3.080) (1.925,3.080) (1.960,3.080) (1.995,3.080) (2.030,3.080) (2.065,3.080) (2.100,3.080) (2.135,3.080) (2.170,3.080) (2.205,3.080) (2.240,3.080) (2.275,3.080) (2.310,3.080) (2.345,3.080) (2.380,3.080) (2.415,3.080) (2.450,3.080) (2.485,3.080) (2.520,3.080) (2.555,3.080) (2.590,3.080) (2.625,3.080) (2.660,3.080) (2.695,3.080) (2.730,3.080) (2.765,3.080) (2.800,3.080) (2.835,3.052) (2.870,2.973) (2.905,2.892) (2.940,2.807) (2.975,2.719) (3.010,2.629) (3.045,2.535) (3.080,2.438) (3.115,2.339) (3.150,2.238) (3.185,2.133) (3.220,2.029) (3.255,1.930) (3.290,1.836) (3.325,1.747) (3.360,1.661) (3.395,1.580) (3.430,1.503) (3.465,1.430) (3.500,1.360) (3.535,1.294) (3.570,1.231) (3.605,1.171) (3.640,1.113) (3.675,1.059) (3.710,1.007) (3.745,0.958) (3.780,0.911) (3.815,0.867) (3.850,0.825) (3.885,0.784) (3.920,0.746) (3.955,0.710) (3.990,0.675) (4.025,0.642) (4.060,0.611) (4.095,0.581) (4.130,0.553) (4.165,0.526) (4.200,0.500) (4.235,0.476) (4.270,0.452) (4.305,0.430) (4.340,0.409) (4.375,0.389) (4.410,0.370) (4.445,0.352) (4.480,0.335) (4.515,0.319) (4.550,0.303) (4.585,0.288) (4.620,0.274) (4.655,0.261) (4.690,0.248) (4.725,0.236) (4.760,0.225) (4.795,0.214) (4.830,0.203) (4.865,0.193) (4.900,0.184) (4.935,0.175) (4.970,0.166) (5.005,0.158) (5.040,0.151) (5.075,0.143) (5.110,0.136) (5.145,0.130) (5.180,0.123) (5.215,0.117) (5.250,0.111) (5.285,0.106) (5.320,0.101) (5.355,0.096) (5.390,0.091) (5.425,0.087) (5.460,0.083) (5.495,0.079) (5.530,0.075) (5.565,0.071) (5.600,0.068) (5.635,0.064) (5.670,0.061) (5.705,0.058) (5.740,0.055) (5.775,0.053) (5.810,0.050) (5.845,0.048) (5.880,0.045) (5.915,0.043) (5.950,0.041) (5.985,0.039) (6.020,0.037) (6.055,0.035) (6.090,0.034) (6.125,0.032) (6.160,0.030) (6.195,0.029) (6.230,0.027) (6.265,0.026) (6.300,0.025) (6.335,0.024) (6.370,0.022) (6.405,0.021) (6.440,0.021) (6.475,0.022) (6.510,0.024) (6.545,0.027) (6.580,0.032) (6.615,0.037) (6.650,0.044) (6.685,0.052) (6.720,0.061) (6.755,0.071) (6.790,0.083) (6.825,0.095) (6.860,0.109) (6.895,0.124) (6.930,0.140) (6.965,0.157) (7.000,0.176) (7.035,0.195) (7.070,0.215) (7.105,0.237) (7.140,0.260) (7.175,0.283) (7.210,0.308) (7.245,0.333) (7.280,0.360) (7.315,0.388) (7.350,0.416) (7.385,0.445) (7.420,0.476) (7.455,0.507) (7.490,0.539) (7.525,0.571) (7.560,0.604) (7.595,0.638) (7.630,0.673) (7.665,0.708) (7.700,0.744) (7.735,0.781) (7.770,0.818) (7.805,0.855) (7.840,0.893) (7.875,0.931) (7.910,0.969) (7.945,1.008) (7.980,1.047) (8.015,1.086) (8.050,1.126) (8.085,1.165) (8.120,1.204) (8.155,1.244) (8.190,1.283) (8.225,1.322) (8.260,1.362) (8.295,1.400) (8.330,1.439) (8.365,1.477) (8.400,1.515) (8.435,1.553) (8.470,1.590) (8.505,1.626) (8.540,1.662) (8.575,1.698) (8.610,1.733) (8.645,1.767) (8.680,1.800) (8.715,1.832) (8.750,1.864) (8.785,1.895) (8.820,1.924) (8.855,1.953) (8.890,1.981) (8.925,2.007) (8.960,2.033) (8.995,2.057) (9.030,2.080) (9.065,2.102) (9.100,2.123) (9.135,2.142) (9.170,2.160) (9.205,2.177) (9.240,2.192) (9.275,2.205) (9.310,2.217) (9.345,2.228) (9.380,2.237) (9.415,2.245) (9.450,2.251) (9.485,2.255) (9.520,2.258) (9.555,2.259) (9.590,2.259) (9.625,2.256) (9.660,2.253) (9.695,2.247) (9.730,2.240) (9.765,2.231) (9.800,2.220) (9.835,2.208) (9.870,2.194) (9.905,2.178) (9.940,2.160) (9.975,2.141) (10.010,2.120) (10.045,2.098) (10.080,2.074) (10.115,2.048) (10.150,2.021) (10.185,1.992) (10.220,1.961) (10.255,1.929) (10.290,1.895) (10.325,1.860) (10.360,1.824) (10.395,1.786) (10.430,1.746) (10.465,1.706) (10.500,1.663)};

\node[font=\small,gray!40!black,anchor=west] at (0.9,3.15) {без торможения: лавина};
\node[font=\small,blue!60!black,anchor=west] at (7.4,1.25) {быстрое торможение};
\node[font=\small,red!70!black,anchor=west] at (7.4,2.55) {медленное торможение};
\end{tikzpicture}
\caption{Частота \nt{ГЛУТ}-группы с обратной связью $w_{EE}=1.5$, $\tau_E=10\ms$, после включения входа. Без торможения (штрих) активность растёт без предела. С торможением силы $w_{EI}w_{IE}=2$ и $\tau_I=2\ms$ (синяя кривая) она выходит на уровень $E^*=h/(1-1.5+2)=0.67h$. С тем же торможением, но медленным, $\tau_I=30\ms$ (красная кривая), активность колеблется.}
\label{fig:ei}
\end{figure}

Считается, что кора работает именно в таком режиме: её собственная обратная связь достаточно сильна, чтобы без торможения уйти в лавину, и держит её в рабочей точке только быстрое торможение~\cite{Tsodyks1997isn}.

У этого режима есть парадоксальная подпись, по которой его можно узнать снаружи~\cite{Tsodyks1997isn}. Добавим в~\eqref{eq:ei} внешний вход $h_I$ прямо \nt{ГАМК}-группе. Установившиеся частоты станут
\begin{equation}
E^* \;=\; \frac{h-w_{EI}\,h_I}{D},\qquad
I^* \;=\; \frac{w_{IE}\,h+(1-w_{EE})\,h_I}{D},\qquad
D \;=\; 1-w_{EE}+w_{EI}w_{IE} .
\label{eq:isn}
\end{equation}
Если $w_{EE}>1$, коэффициент при $h_I$ во второй формуле отрицателен: подтолкнёшь \nt{ГАМК}-нейроны --- и их частота \emph{упадёт}. Причина в петле. Лишнее торможение приглушает \nt{ГЛУТ}-группу, та меньше возбуждает \nt{ГАМК}-нейроны, и они теряют больше, чем получили. При числах рис.~\ref{fig:ei} ($w_{EE}=1.5$, $w_{EI}w_{IE}=2$, $D=1.5$) каждая единица добавленного входа снижает $I^*$ на треть единицы. Эту подпись и нашли в коре: когда ответ нейронов зрительной коры подавляется раздражителем вокруг, тормозящий ток, который они получают, не растёт, а падает вместе с возбуждающим~\cite{Ozeki2009}. Позже ту же подпись нашли в разных областях и слоях коры~\cite{Sanzeni2020}.

Колебания при нарушении условия~(ii) в мозге тоже встречаются. Петля «\nt{ГЛУТ} возбуждает \nt{ГАМК}, \nt{ГАМК} тормозит \nt{ГЛУТ}» с задержкой в несколько миллисекунд порождает ритм с периодом порядка $12$--$50\ms$ (частота $20$--$80$\,Гц), и такие быстрые ритмы в коре хорошо известны~\cite{Cardin2009,Buzsaki2012}. Это не часы компьютера: ритм местный и возникает только тогда, когда сеть активна. Но в течение нескольких циклов он делит время на окна, в которых нейроны могут сработать.

\section{Итог}

\begin{table}[t]
\centering
\small
\begin{tabular}{>{\raggedright}p{3.3cm}>{\raggedright}p{4.4cm}>{\raggedright\arraybackslash}p{4.4cm}}
\toprule
& только \nt{ГЛУТ} & \nt{ГЛУТ} + \nt{ГАМК} \\
\midrule
НЕ & только через датчики «включения---выключения» & запрет $a\wedge\bar b$; НЕ при фоновой активности \\
провода на бит & два & один \\
направление движения & нет & запрет с задержкой \\
регулировка усиления & нет & деление: шунт вместе с фоновым шумом \\
окно совпадения & $\sim\tau$ & сужается торможением \\
выбор одного из многих & нет & взаимное торможение, $\beta>1$ \\
сброс памяти & только усталостью & сигналом через \nt{ГАМК} \\
устойчивость & только усталостью & быстрая отрицательная обратная связь \\
ритм & не нужен & возникает при медленном торможении \\
\bottomrule
\end{tabular}
\caption{Что добавляет \nt{ГАМК} к сети из \nt{ГЛУТ}-нейронов.}
\label{tab:gaba}
\end{table}

\nt{ГАМК} почти не добавляет сети новых \emph{вычислений} --- всё это можно было вычислить и с двухпроводными датчиками, --- а добавляет \emph{управление} (таблица~\ref{tab:gaba}): выбор, сброс, регулировку усиления и устойчивость. Возбуждение в мозге несёт данные, а торможение решает, каким данным пройти, когда и насколько громко. Для каждого четвёртого-пятого нейрона коры это очень разумное разделение труда.

\section{Задачи}

\begin{exercise}[\lvlA]
\label{ex:03:1}
\emph{Шунт в числах.} Постоянная времени мембраны в покое $C/g_L=20\ms$, $E_E=70\mV$.
\par\noindent(а)~По~\eqref{eq:shunt} найдите $V_\infty$ при $g_E=g_L$ и $g_E=4g_L$, без торможения и с $g_I=2g_L$. Сравните с торможением, которое при $g_E=g_L$ вычитает столько же, сколько шунт, и проверьте на втором значении $g_E$, что шунт делит, а не вычитает.
\par\noindent(б)~Найдите $\tau_{\text{эфф}}$ и окно совпадения~\eqref{eq:window} при $\theta=1.5w$ для $g_E=g_L$ без торможения и с $g_I=2g_L$.
\par\noindent(в)~Датчики детектора направления стоят на расстоянии $\Delta x=50\um$, задержка торможения $d=25\ms$. При какой скорости~\eqref{eq:vstar} детектор молчит?
\end{exercise}

\begin{exercise}[\lvlB]
\label{ex:03:2}
\emph{Вывод условий устойчивости.}
\par\noindent(а)~Найдите матрицу линейной системы~\eqref{eq:ei}, её след и определитель, и выведите из них оба условия утверждения~\ref{prop:ei}. Почему при $w_{EE}<1$ условие~(ii) не нужно?
\par\noindent(б)~На границе условия~(ii) при выполненном~(i) собственные числа чисто мнимые. Найдите частоту колебаний на границе и их период при числах рис.~\ref{fig:ei}. Найдите собственные числа при $\tau_I=2\ms$ и $\tau_I=30\ms$ и объясните обе кривые рисунка. Что в модели \texttt{figures/make\_ei.py} удерживает колебания при $\tau_I=30\ms$ от неограниченного роста?
\par\noindent(в)~Выведите~\eqref{eq:isn} и проверьте, что при числах рис.~\ref{fig:ei} $\dd I^*/\dd h_I=-1/3$. При каком условии на $w_{EE}$ подтолкнутые \nt{ГАМК}-нейроны снижают частоту?
\end{exercise}

\begin{exercise}[\lvlB]
\label{ex:03:3}
\emph{Ритм из задержки.} По задаче~\ref{ex:03:2} период колебаний модели~\eqref{eq:ei} выходит в десятки миллисекунд и больше. Заменим медленное торможение запаздывающим: \nt{ГАМК}-посредник отвечает мгновенно, но петля имеет задержку $d$:
\[
\tau_E\frac{\dd E}{\dd t}=-E(t)-G\,E(t-d)+h,\qquad G=w_{EI}w_{IE}>0 .
\]
\par\noindent(а)~Докажите, что при $d=0$ равновесие устойчиво при любом $G$, а при $G<1$ --- при любой задержке.
\par\noindent(б)~При $G>1$ найдите критическую задержку $d_c$, при которой возникают колебания, и их период на пороге. Докажите, что период лежит между $2d_c$ и $4d_c$ и стремится к $4d_c$ при $G\to\infty$.
\par\noindent(в)~Вычислите $d_c$ и период при $\tau_E=10\ms$ для $G=2$ и $G=5$ и сравните с периодами $12$--$50\ms$ быстрых ритмов коры. Проверьте результат, проинтегрировав уравнение методом Эйлера с буфером прошлых значений.
\end{exercise}

\begin{exercise}[\lvlB]
\label{ex:03:4}
\emph{Моделирование: выбор с памятью.} Проинтегрируйте~\eqref{eq:wta} методом Эйлера ($\tau=1$, шаг $10^{-3}$).
\par\noindent(а)~Воспроизведите числа раздела «Выбор»: частоты при $\beta=0.5$ и $\beta=2$ для входов $1.0$ и $0.9$ и сохранение победы после перестановки входов.
\par\noindent(б)~При $\beta=2$ и $I_1=1$ медленно (со скоростью $10^{-3}$ за $\tau$) поднимайте $I_2$ от $0$ до $3$ и опускайте обратно. При каких $I_2$ происходят переключения? Сравните с утверждением~\ref{prop:wta} и найдите ширину петли гистерезиса.
\par\noindent(в)~Стартуя из $r_1=r_2=0$ при $I_1=1+\varepsilon$, $I_2=1$, измерьте время решения (пока проигравший не упадёт ниже $1\,\%$ победителя) при $\varepsilon=10^{-1},\dots,10^{-4}$. Выведите из линейного анализа, на сколько растёт время при уменьшении $\varepsilon$ в десять раз, и сравните с моделированием. Что это значит для выбора между почти равными вариантами?
\end{exercise}

\begin{exercise}[\lvlB]
\label{ex:03:5}
\emph{Проектирование: детектор направления на полосу скоростей.} Схема рис.~\ref{fig:direction}, но детектор должен молчать при движении от $A$ к $B$ с любой скоростью, для которой время пробега $\Delta x/v$ лежит в пределах $5$--$20\ms$. Импульс \nt{ГАМК}-посредника, пришедший к $Y$ через задержку $d_k$ после $A$, запрещает $Y$ срабатывать на отрезке $[d_k,\,d_k+T_I]$, $T_I=5\ms$; возбуждение от $B$ приходит к $Y$ без задержки, и одного импульса $B$ достаточно, чтобы $Y$ сработал.
\par\noindent(а)~Сколько нужно \nt{ГАМК}-посредников и с какими задержками? Докажите, что меньше нельзя.
\par\noindent(б)~Проверьте, что при движении от $B$ к $A$ с любой скоростью $Y$ срабатывает.
\par\noindent(в)~Что делает схема при движении от $A$ к $B$ со скоростью вне полосы? Как изменится ответ (а), если $T_I$ удвоить?
\end{exercise}

\begin{exercise}[\lvlB]
\label{ex:03:6}
\emph{Запреты и фон.}
\par\noindent(а)~Докажите вторую половину утверждения~\ref{prop:gabacomplete}: без фоновой активности сеть из \nt{ГЛУТ}- и \nt{ГАМК}-нейронов вычисляет в точности логические функции с $f(0,\dots,0)=0$. Для полноты постройте схему из одного \nt{ГАМК}-посредника на каждый вход, одного нейрона-запрета на каждый набор входов с $f=1$ и одного нейрона ИЛИ, и сосчитайте нейроны для чётности трёх входов $x\oplus y\oplus z$.
\par\noindent(б)~Пусть \nt{ГАМК}-нейрону разрешено получать входы $x,y,z$ напрямую. Постройте схему чётности трёх входов из двух нейронов: одного \nt{ГАМК} и одного \nt{ГЛУТ}, указав веса и пороги.
\par\noindent(в)~Входы приходят импульсами в разные моменты. Какая гонка сигналов возникает в схеме (б) и как её устранить задержками? Можно ли её устранить, если входы разнесены во времени на неизвестную величину?
\end{exercise}

\begin{exercise}[\lvlC]
\label{ex:03:7}
\emph{Проектирование: сколько посредников нужно для выбора из ста.} Каждая из $n=100$ групп \nt{ГЛУТ}-нейронов должна тормозить каждую другую с силой $\beta$ и не тормозить себя: нужна эффективная матрица связей $W_{ij}=-\beta$ при $i\ne j$, $W_{ii}=0$. \nt{ГАМК}-посредники линейны: частота посредника $m$ равна $\sum_jA_{mj}r_j$, $A\ge0$, и он тормозит группу $i$ с весом $B_{im}\ge0$. Кроме того, группы могут возбуждать друг друга напрямую, $P\ge0$.
\par\noindent(а)~Докажите, что эффективная матрица равна $P-BA$, и что при $P_{ii}\ge0$ одного посредника хватает для любой желаемой $W$.
\par\noindent(б)~Самовозбуждение групп запрещено, $P_{ii}=0$. Докажите, что наименьшее число посредников --- наименьшее $k$ с $\binom{k}{\lfloor k/2\rfloor}\ge n$, и найдите его для $n=100$. Указание: сопоставьте группе $i$ множество посредников, которые её тормозят, и воспользуйтесь теоремой Шпернера о наибольшем семействе подмножеств, ни одно из которых не содержится в другом.
\par\noindent(в)~Прямые возбуждающие связи тоже запрещены, $P=0$. Сколько нужно посредников теперь?
\par\noindent(г)~Сравните три ответа. Какой ценой (задержкой, точностью, числом связей) оплачена экономия в (а) и (б)?
\end{exercise}

\begin{exercise}[\lvlC]
\label{ex:03:8}
\emph{Исследование: когда кора становится сетью, стабилизированной торможением.} Замените в \texttt{figures/make\_ei.py} линейную передаточную функцию \nt{ГЛУТ}-группы на квадратичную $f(u)=[u]_+^2$, оставив \nt{ГАМК}-группу линейной: $\tau_I\,\dd I/\dd t=-I+w_{IE}E+h_I$; параметры рис.~\ref{fig:ei}, $w_{EI}=2$, $w_{IE}=1$. Работайте с копией файла: сам скрипт перезаписывает \texttt{figures/ei\_traces.tex}.
\par\noindent(а)~Эффективная сила обратной связи теперь $w_{EE}f'(u^*)$ и зависит от входа $h$. Найдите аналитически $h_c$, выше которого знак $\dd I^*/\dd h_I$ становится отрицательным, и $h$, при котором нарушается условие~(ii).
\par\noindent(б)~Проверьте (а) моделированием: при $h$ от $0$ до $2$ дайте \nt{ГАМК}-группе добавку $h_I=0.01$ и измерьте изменение $I^*$. Что происходит, если резко включить большой вход $h$ из тишины?
\par\noindent(в)~Открытая часть. Модель предсказывает, что парадоксальный ответ появляется только выше некоторого уровня входа. Предложите опыт, который проверяет это предсказание, и обсудите, что следует из того, что подпись~\eqref{eq:isn} находят в коре и при слабой, фоновой активности. Какие свойства модели надо изменить, чтобы она это воспроизвела?
\end{exercise}
